\subsection{Geometric realization of the weight-lowering obstruction}

Retain the notation of the Hasse-duality subsection.  In particular,
$P=p(p-1)$, $S=\operatorname{div}(A)$, $Z=pS$, and
$Q_r=H^0(Z,\omega^r\vert_Z)$.  Let $\widetilde A=E_{p-1}$ be the
normalized integral lift of the Hasse invariant.  Its $q$-expansion is
congruent to $1$ modulo $p$, and therefore
\begin{equation}\label{eq:hasse-p-lift}
 \widetilde A^p\equiv1\pmod {p^2}
\end{equation}
as a $q$-expansion.

\begin{proposition}[The Bockstein class is a first Hasse jet]
\label{prop:geometric-bockstein}
Let $g$ be a divided congruence modulo $p^2$.  Suppose that
\begin{enumerate}
\item $g$ has a representative $G\in M_{w+P}(\mathbb Z_{(p)})$ modulo
      $p^2$; and
\item $g\bmod p$ has a representative $h\in M_w(\mathbb Z_{(p)})$.
\end{enumerate}
Then
\[
 \Gamma_w(g):=
 \left.\left(\frac{G-\widetilde A^p h}{p}\bmod p\right)\right|_Z
 \in Q_{w+P}
\]
is independent of all choices.  Moreover,
\[
 \Gamma_w(g)=0
 \quad\Longleftrightarrow\quad
 g\text{ has a representative in }M_w(\mathbb Z_{(p)})
       \text{ modulo }p^2.
\]
Under the $q$-expansion identification, $\Gamma_w(g)$ is the geometric
image of the divided-congruence class $\beta_w(g)$ from
Lemma~\ref{lem:bockstein}.
\end{proposition}

\begin{proof}
The numerator is divisible by $p$: both terms have weight $w+P$, and their
$q$-expansions agree modulo $p$ because $\widetilde A^p\equiv1\pmod p$.
Thus the displayed quotient reduces to a form of weight $w+P$ modulo $p$.

Changing $G$ without changing its $q$-expansion modulo $p^2$ changes the
quotient by zero modulo $p$, by the saturated $q$-expansion principle.
If $h'$ is another choice, the integral echelon argument used in
Lemma~\ref{lem:bockstein} gives $h'-h=pm$ for some $m\in M_w$.  The two
quotients then differ modulo $p$ by
\[
 -A^p m\in A^pM_w.
\]
Its restriction to $Z$ is zero.  This proves independence.

The restriction sequence for the Cartier divisor $Z$ is
\[
 0\longrightarrow\omega^w
 \xrightarrow{\ A^p\ }\omega^{w+P}
 \longrightarrow\omega^{w+P}\vert_Z\longrightarrow0.
\]
Hence $\Gamma_w(g)=0$ exactly when
\[
 \frac{G-\widetilde A^p h}{p}\equiv A^pm\pmod p
\]
for some $m\in M_w$.  In that case
\[
 G\equiv\widetilde A^p(h+pm)\pmod {p^2}.
\]
By \eqref{eq:hasse-p-lift}, the $q$-expansion of the right side is that of
$h+pm$ modulo $p^2$, so $h+pm$ is a weight-$w$ representative of $g$.
The converse follows by reversing the argument.  Finally, before
restriction to $Z$ the displayed quotient has the same $q$-expansion as
$(g-h)/p$ modulo $p$; quotienting by $A^pM_w$ is precisely its realization
in the first-Hasse-jet space.  This proves the last assertion.
\end{proof}

The index shift in Proposition~\ref{prop:geometric-bockstein} is essential:
the obstruction to lowering to weight $w$ belongs to $Q_{w+P}$, not to
$Q_w$.

\begin{lemma}[Theta commutes with the integral Hasse inclusion]
\label{lem:theta-hasse-inclusion}
For every integral $q$-series $h$ and every $m\geq 0$,
\[
 \theta^m(\widetilde A^p h)
 \equiv \widetilde A^p\theta^m(h)\pmod {p^2}.
\]
\end{lemma}

\begin{proof}
Write $\widetilde A=1+pa(q)$.  Then $\theta\widetilde A$ is divisible by
$p$, and hence
\[
 \theta(\widetilde A^p)
 =p\widetilde A^{p-1}\theta\widetilde A
 \equiv0\pmod {p^2}.
\]
The product rule gives the assertion for $m=1$, and iteration proves it
for every $m$.
\end{proof}

\begin{lemma}[A horizontal unit of the required weight]
\label{lem:horizontal-unit}
After pullback to a compactified modular curve with auxiliary full level
$3$, there is a form
\[
 F\in H^0(X,\omega^{3p-1})
\]
whose restriction to $S$ is nowhere zero.  Consequently
\[
 U:=F^p\vert_Z\in H^0(Z,\omega^{p(3p-1)}\vert_Z)
\]
is a unit and $\Theta(U)=0$.
\end{lemma}

\begin{proof}
Multiplication by $A$ gives
\[
 0\longrightarrow\omega^{2p}
 \xrightarrow{\ A\ }\omega^{3p-1}
 \longrightarrow\omega^{3p-1}\vert_S\longrightarrow0.
\]
Serre duality and Kodaira--Spencer give
\[
 H^1(X,\omega^{2p})^\vee
 \simeq H^0(X,\omega^{2-2p}(-C))=0.
\]
The restriction map onto $S$ is therefore surjective.  Since $S$ is a
finite reduced scheme, the restricted line bundle has a global generator;
lift such a generator to $F$.  It remains a unit on the nilpotent
thickening $Z$.  Its $p$-th power is a unit of weight $p(3p-1)$, and the
product rule in characteristic $p$ gives $\Theta(F^p)=0$.
\end{proof}

Now put
\[
 w_-=k+2i-P,
 \qquad
 w_+=k+2i^\dagger+P,
\]
and define the actual geometric Bockstein weights
\[
 r_-:=w_-+P=k+2i,
 \qquad
 r_+:=w_++P=k+2i^\dagger+2P.
\]
They satisfy
\begin{equation}\label{eq:natural-bockstein-weight-sum}
 r_-+r_+=P+2+p(3p-1).
\end{equation}

\begin{theorem}[Auxiliary-level natural Bockstein pairing]
\label{thm:natural-bockstein-pairing}
For a choice of $F$ as in Lemma~\ref{lem:horizontal-unit}, the formula
\[
 [x,z]^{\mathrm{nat}}:=\langle x,U^{-1}z\rangle_{r_-}
\]
defines a perfect pairing
\[
 Q_{r_-}\times Q_{r_+}\longrightarrow\mathbb F_p.
\]
Theta is skew-adjoint under this pairing.  The construction is not
canonical in $F$, but every such choice is perfect and detects the zero
class.  Pullback to auxiliary level is faithful, so vanishing of the
level-one Bockstein is equivalent to vanishing of its pulled-back class.
\end{theorem}

\begin{proof}
Equation~\eqref{eq:natural-bockstein-weight-sum} and the invertibility of
$U$ identify the second factor with $Q_{P+2-r_-}$.  Perfectness follows
from Proposition~\ref{prop:hasse-perfect-pairing}.  Since
$\Theta(U)=\Theta(U^{-1})=0$, Proposition~\ref{prop:theta-adjointness}
gives
\[
 [\Theta x,z]^{\mathrm{nat}}
 =-[x,\Theta z]^{\mathrm{nat}}.
\]
Faithful pullback proves the final assertion.
\end{proof}

\begin{remark}
The perfect pairing in Theorem~\ref{thm:natural-bockstein-pairing} depends
on the choice of the auxiliary-level unit $F$.  The theorem uses it only to
identify dual vector spaces and detect the zero class after faithful
pullback.  It does not assert a canonical level-one bilinear form or descend
the chosen unit through the tame deck group.  By contrast, the unshifted
Hasse-quotient pairing in Proposition~\ref{prop:hasse-perfect-pairing} is
canonical and does descend by invariants.
\end{remark}

For the complementary theta iterates, Proposition
\ref{prop:geometric-bockstein} now gives the correctly typed classes
\[
 \Gamma_{w_-}(\theta^if)\in Q_{r_-},
 \qquad
 \Gamma_{w_+}(\theta^{i^\dagger}f)\in Q_{r_+}.
\]
The remaining cross-wedge problem is no longer the construction of their
ambient pairing.  It is the following concrete compatibility statement:
under the isomorphism $Q_{r_+}\simeq Q_{r_-}^\vee$ induced by
$[\ ,\ ]^{\mathrm{nat}}$, the second displayed class must equal the
functional obtained from the first by the integral theta-iterate
construction, up to the explicit nonzero scalar in the $D=2$ block
formula.  Perfectness and theta adjointness alone do not prove that
identification.

Lemma~\ref{lem:theta-hasse-inclusion} proves the Hasse-inclusion part of
the required square.  Theorem~\ref{thm:explicit-periodic-first-jet} gives
an explicit formula for the lower minimal-weight projection.  The
unresolved part is its factorial block reduction and the corresponding
compatibility with the complementary class.
